\documentclass{article}

\usepackage[preprint]{neurips_2026}

\usepackage[utf8]{inputenc}
\usepackage[T1]{fontenc}
\usepackage{hyperref}
\usepackage{url}
\usepackage{booktabs}
\usepackage{graphicx}
\usepackage{amsfonts}
\usepackage{nicefrac}
\usepackage{microtype}
\usepackage{xcolor}
\usepackage{enumitem}
\usepackage{fancyvrb}
\usepackage{listings}
\lstdefinestyle{kisspy}{
  language=Python,
  basicstyle=\ttfamily\scriptsize,
  keywordstyle=\color[rgb]{0.0,0.0,0.8}\bfseries,
  commentstyle=\color[rgb]{0.25,0.5,0.35}\itshape,
  stringstyle=\color[rgb]{0.73,0.13,0.13},
  showstringspaces=false,
  frame=single,
  framesep=2mm,
  numbers=left,
  numberstyle=\tiny\color{gray},
  numbersep=5pt,
  breaklines=true,
  tabsize=4,
  captionpos=b,
}
\lstnewenvironment{prompt}{%
  \lstset{language={},basicstyle=\ttfamily\scriptsize,frame=single,framesep=2mm,breaklines=true,columns=fullflexible}%
}{}

\title{KISS Sorcar: A Stupidly-Simple General-Purpose and Software Engineering AI Assistant}

\author{%
Koushik Sen \\ EECS Department, UC Berkeley\\ \texttt{ksen@berkeley.edu} \\ }

\begin{document}

\maketitle

\begin{center}
\textit{"Everything should be made as simple as possible, but not simpler." — Albert Einstein}
\end{center}


\begin{abstract}
Large language models can generate code and call tools with remarkable fluency, yet deploying them as practical software engineering assistants still exposes stubborn gaps: finite context windows, a single mistake that can derail entire sessions, agents that get stuck in dead ends, AI slop, and generated changes that are difficult to review or revert.

We present \textbf{KISS Sorcar}, a general-purpose assistant and integrated development environment (IDE) built on top of the \textbf{KISS Agent Framework}, a stupidly-simple AI agent framework of roughly 1,900~lines of code.  The framework addresses these gaps using a robust system prompt and through a five-layer agent hierarchy in which each layer adds exactly one concern: budget-tracked ReAct execution, automatic continuation across sub-sessions via summarization, coding, and browser tools with parallel sub-agents, persistent multi-turn chat with history recall, and git worktree isolation so every task runs on its own branch.  Both KISS Sorcar and the KISS Agent Framework are grounded in disciplined software engineering practice; these principles are encoded directly into the agent's system prompt, enabling KISS Sorcar to write code that is simple, elegant, maintainable, and bug-free.

To assess the power of the KISS agent framework, we implemented KISS Sorcar as a free, open-source Visual Studio Code extension that runs locally, is effective for long-horizon tasks, and supports browser automation, multimodal input, and Docker containers.

In this research, we deliberately prioritize output quality over speed:  giving a frontier model adequate time to validate its own output---running linters, type checkers, and tests---dramatically reduces the low-quality code that plagues faster but less thorough agents.  The entire system was built using itself in 4 months, providing a continuous stress test in which any bug was patched immediately after its manifestation.  On Terminal~Bench~2.0, KISS Sorcar achieves a 62.2\% overall pass rate with Claude Opus~4.6, comparing favorably to Claude~Code (58\%) and Cursor Composer~2 (61.7\%).  We believe that our results are significant because we didn't tune our prompts or any model specifically for the Terminal Bench 2.0.
\end{abstract}

%----------------------------------------------------------------------
\section{Introduction}
\label{sec:intro}
%----------------------------------------------------------------------

Modern Large language models (LLMs), such as Anthropic's Claude Opus 4.7~\citep{anthropic2026opus47}, OpenAI's GPT 5.5~\citep{openai2026gpt55}, and Google's Gemini 3.1~\citep{google2026gemini31}, have demonstrated a remarkable ability to generate code, reason about software architecture, and use developer tools~\citep{chen2021evaluating,roziere2023code}.  A growing body of work has explored how to harness these capabilities for autonomous software engineering, from single-session agents that resolve GitHub issues~\citep{yang2024sweagent,wang2024openhands} to industrial products marketed as AI software and general assistants ~\citep{copilot2021,cursor2024,devin2024,claudecode2025,codex2025,openclaw2025}. Yet using an LLM as a practical software engineering or general assistant still exposes several stubborn gaps: context windows are finite, a single mistake can derail an entire session, agents get stuck in dead ends, models generate AI slop, and generated changes are difficult to review or revert once applied to a live codebase.

We propose the \emph{KISS Agent Framework}, a stupidly simple AI agent framework containing around 1,900 lines of code for the core agent implementation.  We try to address the above-mentioned gaps through a robust system prompt (which we describe in depth in Section~\ref{sec:system-prompt}) and a five-layer agent hierarchy in which each layer solves exactly one concern:

\begin{enumerate}[nosep]
  \item \textbf{KISS Agent} --- budget-tracked
ReAct~\citep{yao2023react} loop with native function calling.
  \item \textbf{Relentless Agent} --- automatic summarization and continuation 
across sub-sessions.
  \item \textbf{Sorcar Agent} --- coding tools, browser automation,
and parallel sub-agent execution.
  \item \textbf{Chat Sorcar Agent} --- persistent multi-turn chat
sessions with history recall.
  \item \textbf{Worktree Sorcar Agent} --- git worktree isolation so
every task runs on its own branch.
\end{enumerate}

The name ``KISS'' reflects the \emph{Keep It Simple, Stupid} design philosophy in software engineering: each layer is small, each concern is isolated, and the overall system avoids unnecessary abstraction.

\begin{figure}[t]
\centering
\includegraphics[width=\textwidth]{KISS-Sorcar.png}
\caption{Screenshot of KISS Sorcar running as a VS~Code extension. The sidebar shows the agent's chat interface with real-time budget tracking, while the editor displays the code being modified.}
\label{fig:kiss-sorcar-screenshot}
\end{figure}

To access and enhance the power and capabilities of the KISS agent framework, we implement KISS Sorcar, a general-purpose assistant and an integrated development environment (IDE) built on top of it. It codes really well and works pretty fast.  We implemented KISS Sorcar as a Visual Studio Code extension and runs locally. It has full browser support (using open-source Chromium browser and Playwright), multimodal support, Docker container support, OpenClaw-like features (whose discussion is beyond the scope of the paper), and a mobile/web app. The good part is that KISS Sorcar is completely free and open-source; all one needs is a model API key from a major LLM provider, such as Anthropic.  We implemented this framework in roughly 4 months, and the repository is available at \url{https://github.com/ksenxx/kiss_ai}. The name ``Sorcar'' pays homage to P.~C.~Sorcar, the legendary Bengali magician, evoking the idea of an agent that performs feats that appear magical yet are grounded in disciplined engineering.  Indeed, both KISS Sorcar and the KISS Agent Framework are grounded in disciplined engineering practice: these principles are encoded directly into the agent's system prompt, enabling KISS Sorcar to write code that is simple, elegant, maintainable, and bug-free.

We want to emphasize that KISS Sorcar has been built using itself.  The entire codebase---the KISS Agent framework, the Sorcar agent layers, the VS~Code extension, and the system prompt---was developed by KISS Sorcar operating on its own repository.  This self-hosting discipline provides a continuous-integration-style stress test: if the agent introduces a bug that impairs its ability to function, we immediately ask the agent to fix it by analyzing the trajectory and code of KISS Sorcar.  The simplicity of the layered architecture was both a prerequisite for and a consequence of this bootstrapping process. The simplicity of the framework enabled the agent to implement features confidently without bugs and AI slop.  The five core agent classes are remarkably compact: the KISS Agent comprises 413~lines, the Relentless Agent 321~lines, the Sorcar Agent 322~lines, the Chat Sorcar Agent 125~lines, and the Worktree Sorcar Agent 704~lines---a total of roughly 1,885~lines of code (excluding empty lines and comments).

In the project, we deliberately prioritize output quality over speed.  In our experience, using a weaker or cheaper model often forces the developer to discard the agent's work and retry, ultimately increasing the total cost of completing a task.  Conversely, giving a frontier model adequate time to validate its own output---running linters, type checkers, and tests before declaring success---dramatically reduces the ``slop'' (low-quality, subtly incorrect code) that plagues faster but less thorough agents.  We expect token costs and inference latencies to continue to fall~\citep {gao2025trae}, making this quality-first posture increasingly practical.  In the meantime, the code produced by our agent is consistently well-organized, simple, and idiomatic.

We evaluate on Terminal Bench~2.0 and achieve a 62.2\% overall pass rate using Claude Opus~4.6, comparing favorably to Claude~Code (58\%) and Cursor Composer~2 (61.7\%)~\citep{cursor2026composer2} on the same benchmark (Section~\ref{sec:evaluation}).  We believe that our results are significant because we didn't tune our prompts or any model specifically for the Terminal Bench 2.0.

In KISS Sorcar, we deliberately kept the system simple.  We included only the agent technologies necessary for KISS Sorcar to function as a robust general-purpose software engineering assistant. We showed that a simple agent framework without sophisticated agent technologies, such as trajectory compaction and asynchronous multi-agent orchestration, was not necessary to build KISS Sorcar.  By building KISS Sorcar using itself and beating both the Cursor and Claude Code agents, we demonstrated that time-tested software engineering techniques and principles are invaluable for building reliable, practical agent systems.

\paragraph{Outline.}
Section~\ref{sec:architecture} presents the five-layer agent architecture and its motivating design principles. Section~\ref{sec:evaluation} reports evaluation results on Terminal Bench~2.0.  Section~\ref{sec:system-prompt} details the system prompt. Section~\ref{sec:vscode-features} covers the VS~Code extension. Section~\ref{sec:painless} illustrates painless software engineering through a real development session. Section~\ref{sec:related} discusses related work, and Section~\ref{sec:conclusion} concludes.

%----------------------------------------------------------------------
\section{Agent Architecture}
\label{sec:architecture}
%----------------------------------------------------------------------

We initially built the KISS Agent Framework to rapidly prototype and experiment with various prompt optimization techniques, such as Gepa~\citep{agrawal2025gepa}, and evolutionary algorithms for algorithmic and code optimization, such as AlphaEvolve~\citep{novikov2025alphaevolve} and OpenEvolve~\citep{openevolve2025}. We focused heavily on keeping the agent framework simple so we could rapidly prototype and experiment with ideas.  The simplicity of the framework also enabled coding agents to write simple, bug-free code.  We ultimately did not use any prompt optimization techniques when creating the system prompt for KISS Sorcar, as we hand-tuned it based on our long-term experience with KISS Sorcar and its behavior.  The KISS Agent Framework uses five agent layers in a layered architecture combining composition and inheritance. Each layer delegates upward for the concerns it does not own.

%......................................................................
\subsection{KISS Agent}
\label{sec:kiss-agent}
%......................................................................

The KISS Agent is the innermost execution unit.  It implements a standard ReAct loop~\citep{yao2023react} with the following characteristics.  

\begin{lstlisting}[style=kisspy,numbers=none,caption={A complete KISS agent with a single tool.},label={lst:simple-agent}]
from kiss.core.kiss_agent import KISSAgent

def calculate(expression: str) -> str:
    """Evaluate a math expression."""
    return str(eval(expression))

agent = KISSAgent(name="Math Buddy")
result = agent.run(
    model_name="gemini-2.5-flash",
    prompt_template="Calculate: {question}",
    arguments={"question": "What is 15% of 847?"},
    tools=[calculate]
)
print(result)  # 127.05
\end{lstlisting}

\textbf{Native function calling.}  We register tools as ordinary Python callables.  The agent builds an OpenAI-compatible tool schema once at setup time and caches it, avoiding redundant schema construction on every LLM call.  A special \texttt{finish} tool signals task completion and returns the result to the caller.

\textbf{Step, token, and budget tracking.}  At every step, the agent extracts input and output token counts from the API response, computes the dollar cost using a per-model pricing table, and updates both a local budget counter and a global (cross-agent) budget counter protected by a class-level lock.  The agent checks three limits before each step: the per-agent budget, the global budget, and the maximum step count.

\textbf{Error resilience.}  The agent retries transient API errors (rate limits, server errors) up to a configurable threshold of consecutive failures.  It detects non-retryable errors (authentication failures, permission denials) and raises them immediately.

\textbf{Non-agentic mode.}  When tools are not needed, the agent can run a single generation without the ReAct loop, which is useful for summarization or question-answering sub-tasks.

Listing~\ref{lst:simple-agent} shows a complete, working agent in under ten lines of code.  The developer defines an ordinary Python function (\texttt{calculate}), instantiates a \texttt{KISSAgent}, and calls its \texttt{run} method with a model name, a prompt template, template arguments, and a list of tools.  The framework automatically handles tool-schema generation, the ReAct loop, and budget tracking.

The KISS Agent is stateless across runs: each call to its run method resets the conversation, token counters, and tool registry.  This makes it safe to reuse a single agent instance for multiple sequential tasks.

%......................................................................
\subsection{Relentless Agent}
\label{sec:relentless-agent}
%......................................................................

The Relentless Agent wraps a KISS Agent in a continuation loop.  Its core contribution is the ability to execute tasks that exceed a single context window by breaking them into sub-sessions.

Rather than investing in context-compaction techniques, we adopt a simple continuation protocol: when a sub-session exhausts its context window or step budget, the agent produces a \emph{structured summary} of every action taken so far---chronologically ordered, with explanations and relevant code snippets---and a fresh sub-session resumes from that summary.  This approach is related in spirit to Reflexion~\citep{shinn2023reflexion}, which feeds verbal self-critiques back into subsequent trials, but uses Reflexion-like technique to continue a task.  While developing KISS Sorcar, we found in our experience that a na\"ive instruction to ``summarize the current context'' produced poor continuations; requiring a \emph{step-by-step chronological account with code snippets} dramatically improved coherence across sub-sessions. A potential limitation is that summaries may grow unwieldy for multi-day tasks; in practice, we have not encountered this problem even for tasks spanning several hours, but a thorough evaluation of summary scaling remains future work.

\textbf{Continuation protocol.}  The \texttt{finish} tool exposed to the inner KISS Agent accepts three fields: a success flag, a continue flag, and a summary.  When the agent sets \texttt{is\_continue=True}, the Relentless Agent starts a new sub-session with a fresh context window.  The prompt for the new session includes a chronologically ordered list of all prior attempt summaries and instructs the agent not to redo completed work.  The continuation prompt template is:

\begin{prompt}
# Task Progress (Continuation {continuation_number})

{progress_text}

# Continue
- Complete the rest of the task.
- **DON'T** redo completed work.
- If you have been retrying the same approach without progress,
  step back and rethink the strategy from scratch.
\end{prompt}

\textbf{Forced continuation on failure.}  If a sub-session raises an exception (e.g., the step limit is hit before the agent calls finish), the Relentless Agent does not abort.  Instead, it saves the full trajectory to a temporary file, spawns a separate summarizer agent to read the trajectory and produce a concise summary, and then uses that summary as the progress text for the next sub-session.  This ensures that even crashed sessions contribute useful context to subsequent attempts.  The summarizer receives the following prompt:

\begin{prompt}
# Summarizer

The trajectory of the agent is stored in the file: {trajectory_file}

# Instructions
- Read the trajectory file and analyze it.  The trajectory file
  could be large.
- Return a precise chronologically-ordered list of things the
  agent did, with the reason for doing that, along with relevant
  code snippets
\end{prompt}

\noindent To force the agent to self-continue before hitting the step limit, we augment the system prompt with an instruction that fires near the end of the budget:

\begin{prompt}
# MOST IMPORTANT INSTRUCTIONS
- **At step {step_threshold}: you MUST call
  finish(success=False, is_continue=True,
  summary="precise chronologically-ordered list of things
  the agent did with the reason for doing that along with
  relevant code snippets")** or if the task is not complete
  and you are at risk of running out of steps or context
  length.
- Work dir: {work_dir}
- Current process PID: {current_pid} -- NEVER kill this
  process.
\end{prompt}

%......................................................................
\subsection{Sorcar Agent}
\label{sec:sorcar-agent}
%......................................................................

The Sorcar Agent adds the tools that make the system useful for software development and general-purpose automation.

\textbf{Coding tools.}  We provide four core tools: a shell command executor with streaming output, a file reader, a precise string-based file editor, and a file writer.  The shell executor supports a configurable timeout, streams output to the user interface in real time, and respects a stop event that allows the user to cancel a running command.  Note that we kept the tool names (Bash, Edit, Write, Read) the same as in Claude Code because the underlying Anthropic models do not make mistakes with these tool names.

\textbf{Browser automation.}  A web-use tool provides programmatic browser control: navigating to URLs, reading page accessibility trees, clicking elements, typing text, pressing keys, scrolling, and taking screenshots.  This enables the agent to research documentation, verify deployed applications, and interact with web-based tools. It uses the open-source Chromium browser using the Playwright library. 

\textbf{Parallel sub-agents.}  An optional parallel execution tool spawns independent Sorcar Agent instances in a thread pool.  Each sub-agent gets its own LLM context and tool set.  This is useful for embarrassingly parallel tasks such as summarizing multiple files or researching independent topics.  We collect the results and return them in input order.  We do not activate this tool by default because in the IDE, we cannot stream multiple agent outputs coherently in the chat window.

\textbf{User interaction.}  An ask-user-question tool allows the agent to pause execution and request clarification from the user.  In the VS~Code integration, this renders as a text input in the sidebar; in CLI mode, it reads from standard input.

\textbf{Docker isolation.}  When a Docker image is specified, we replace the coding tools with Docker-aware variants that execute commands inside a container, providing an additional layer of sandboxing for untrusted tasks.

%......................................................................
\subsection{Chat Sorcar Agent}
\label{sec:chat-agent}
%......................................................................

The Chat Sorcar Agent adds multi-turn conversation persistence.

\textbf{Chat sessions.}  We assign each task to a chat session identified by a stable chat ID.  The agent persists tasks and their results to a local database (\texttt{sorcar.db}).  When a new task arrives within the same chat session, the agent loads prior tasks and results and prepends them to the prompt as numbered context entries, allowing the LLM to reference earlier work.

\textbf{Bounded chat context.}  To prevent unbounded growth of the prompt as a chat session accumulates many tasks, the agent caps the number of in-context history entries at \texttt{MAX\_TASKS=10}.  When the cap is exceeded, it preserves the first two entries (which typically establish the user's overall intent for the session) and the most recent entries, dropping the middle entries that are least likely to be referenced.  This keeps the context relevant and bounded while retaining both the session's framing and its current state.

\textbf{Session management.}  The agent supports three operations: starting a new chat (with a fresh chat ID), resuming a chat by task description (which looks up the corresponding chat ID), and resuming by explicit chat ID.  This enables both automatic session continuity in an IDE and manual session selection from the command line.

\textbf{Frequent task tracking.}  Each time a task is executed, the agent records the task description in a frequency table.  The IDE sidebar surfaces the most frequent tasks so users can re-issue them with one click, turning recurring requests (``run the test suite,'' ``regenerate the changelog'') into a click-to-replay experience.

\textbf{Metadata persistence.}  After each task, the agent records metadata including the model used, working directory, software version, token count, cost, and whether the task used parallel execution or worktree isolation.  This audit trail supports cost analysis and debugging.

%......................................................................
\subsection{Worktree Sorcar Agent}
\label{sec:worktree-agent}
%......................................................................

The Worktree Sorcar Agent is the outermost layer and the one that users interact with in the VS~Code extension and the Sorcar web app.  Its defining feature is git-worktree isolation.

\textbf{Branch-per-task.}  When a task starts, the agent creates a new git branch and a corresponding worktree directory.  The branch name encodes the chat ID and a timestamp for uniqueness.  All agent modifications happen inside the worktree; the user's main working tree remains untouched.

\textbf{Dirty-state preservation.}  If the user's main working tree has uncommitted changes, the agent copies them into the worktree and creates a baseline commit.  This ensures the agent sees the same state as the user, while keeping the user's actual index and working tree clean.  During merge, we use cherry-pick from the baseline commit to replay only the agent's changes, excluding the dirty-state snapshot.

\textbf{Concurrency safety.}  A per-repository file lock serializes the checkout, stash, merge, and pop sequence so that concurrent tabs in the IDE cannot interleave operations on the same repository. Thread-local storage isolates per-task state (stream parsing buffers, bash output buffers, recording state) so that stopping one task does not corrupt another.

\textbf{Crash recovery.}  We store all worktree state in git itself (branch names, git config entries) rather than in sidecar files.  On process restart, the agent queries git for any pending branch matching its chat ID prefix and reconstructs all instance attributes from git config, enabling seamless recovery.

\textbf{Graceful fallback.}  If the working directory is not inside a git repository, if the repository has no commits, or if HEAD is detached, the agent falls back to direct execution without worktree isolation, ensuring it never fails due to git preconditions.


%----------------------------------------------------------------------
\section{Evaluation on Terminal Bench~2.0}
\label{sec:evaluation}
%----------------------------------------------------------------------

Before we discuss a crucial aspect of KISS Sorcar, the system prompt, in a lengthy section, we describe the evaluation outcome of KISS Sorcar on the Terminal Bench 2.0, which was also recently used by the Cursor agent of Composer 2.0.

We evaluate our system on Terminal Bench~2.0,\footnote{\url{https://www.tbench.ai/}} a benchmark comprising 89 diverse terminal-based programming tasks, ranging from building legacy compilers and configuring servers to solving cryptanalysis challenges and training machine-learning models.  Each task runs in an isolated Docker container; a separate verifier automatically judges the result.  We use the Harbor\footnote{\url{https://github.com/harbor-framework/harbor}} framework to orchestrate execution, and Claude Opus~4.6 as the underlying LLM.  We do not modify the general system prompt or inject Terminal Bench 2.0-specific instructions during the evaluation.  We carried out our evaluation on a 2025 MacBook Air 15" with an M4 processor and 24GB RAM.

\subsection{Setup}

We run 5 independent trials per task.  The agent is \texttt{SorcarHarborAgent}, a thin Harbor adapter that installs and invokes the Sorcar CLI inside each container.  We hard-skip 9 tasks that we verified to be infeasible for Opus~4.6 across 6+ prior attempts (e.g.\ CompCert compilation, Windows~3.11 GUI installation, video OCR) to save time and token cost.  Skipped tasks still count as failures.

\subsection{Aggregate Results}

Table~\ref{tab:tb2-aggregate} summarizes the aggregate statistics.

\begin{table}[h]
\centering
\caption{Terminal Bench~2.0 aggregate results (89 tasks, 5 trials each, Claude Opus~4.6).}
\label{tab:tb2-aggregate}
\begin{tabular}{lr}
\toprule
\textbf{Metric} & \textbf{Value} \\
\midrule
Total tasks              & 89 \\
Overall pass rate        & 62.2\% (277/445) \\
pass@any (at least 1/5 passes) & 78.7\% (70/89) \\
pass@all (all 5 pass)    & 43.8\% (39/89) \\
Always-fail tasks        & 19 \\
Always-pass tasks        & 39 \\
Mixed-result tasks       & 31 \\
Median cost per trial    & \$0.45 \\
Mean cost per trial      & \$0.90 \\
Median duration per trial & 202\,s \\
Mean duration per trial  & 446\,s \\
\bottomrule
\end{tabular}
\end{table}

The 62.2\% overall pass rate compares favorably to other agents using the same underlying model: at the time of writing, Claude~Code (also Opus~4.6) scores approximately 58\% on the Terminal Bench~2.0 leaderboard, and Cursor's Composer~2---a custom fine-tuned model trained with large-scale reinforcement learning~\citep{cursor2026composer2}---achieves 61.7\%.  Our result suggests that the layered architecture and the structured system prompt described in Sections~\ref{sec:architecture}--\ref{sec:system-prompt} contribute meaningfully beyond what the base model provides.

\subsection{Task-Level Breakdown}

\textbf{Consistently solved tasks (39 of 89).} These include cryptanalysis (FEAL differential), game-playing (chess best move), git operations (leak recovery), server configuration (gRPC key-value store, PyPI server, NGINX logging), data processing (resharding), formal verification (Coq~\texttt{plus\_comm}), ML inference (HuggingFace model serving, LLM batching scheduler), and system emulation (QEMU startup).  The breadth of this set---spanning systems, security, data engineering, and formal methods---demonstrates that the agent's generality is not limited to a single domain.

\textbf{Consistently failed tasks (19 of 89).} The failures cluster into three categories: (1)~tasks requiring graphical or multimedia capabilities unavailable in the container (video processing, Windows~3.11 GUI, MTEB leaderboard scraping, extracting moves from video), (2)~tasks demanding very long or resource-intensive builds that exceed the container's time or memory limits (CompCert, Doom for MIPS, Caffe CIFAR-10, training fastText on Yelp data), and (3)~tasks with niche domain-specific requirements that the model struggles to satisfy (DNA insertion, OCaml GC patching, polyglot C/Python binaries, protein assembly, cell segmentation).

\textbf{Mixed-result tasks (31 of 89).} Tasks such as \texttt{write-compressor} (3/5), \texttt{crack-7z-hash} (4/5), and \texttt{feal-linear-cryptanalysis} (4/5) succeed in most trials but occasionally fail due to non-determinism in the model's reasoning or timing-sensitive environment interactions.  Conversely, \texttt{cancel-async-tasks} (1/5) and \texttt{dna-assembly} (1/5) succeed rarely, suggesting they are at the boundary of the model's capability.

\textbf{Leaderboard context.}  KISS Sorcar does not score as high as other coding agents reported at the Terminal Bench 2.0 leaderboard, but we believe that our results are significant because we didn't tune our prompts or any model specifically for the Terminal Bench 2.0.  We simply used the general system prompt and the Claude Opus 4.6 model.  Regarding low score compared to other coding agents, we wanted to note that recent analysis has found widespread cheating on popular agent benchmarks, including Terminal Bench~2.0: the top three submissions commit harness-level cheating (e.g.\ leaking verifier code or answer keys into the agent's environment), and task-level cheating (e.g.\ Googling answers, mining git history, hardcoding test outputs) affects 28+ submissions across 9 benchmarks~\citep{stein2026cheatingblog,stein2026detecting}. Separately, we discovered, using an automated benchmark audit agent, that 45 confirmed hacking solutions across 13 widely used benchmarks exhibited process-isolation failures, answer leakage, and weak test assertions that allow perfect scores without solving a single problem~\citep{wang2026trustworthyblog}.

%----------------------------------------------------------------------
\section{The System Prompt}
\label{sec:system-prompt}
%----------------------------------------------------------------------

The system prompt is a structured document that governs the agent's behavior across all tasks.  It is not a generic instruction to ``be helpful'' but rather a precise specification of an engineering discipline. We describe its key sections.

%......................................................................
\subsection{Execution Mindset}
%......................................................................

The prompt opens with a directive that sets the tone:

\begin{prompt}
# FOCUS ON THE GIVEN TASK. ITS COMPLETION IS YOUR SOLE GOAL.
# BE RELENTLESS. BE CALM. BE RIGOROUS. BE ACCURATE. CHECK FACTS. NO AI SLOP.
\end{prompt}

\noindent Each directive addresses a specific failure mode observed in
LLM-based agents:

\textbf{``FOCUS ON THE GIVEN TASK. ITS COMPLETION IS YOUR SOLE GOAL.''} LLMs have a tendency to drift: they comment on the difficulty of a problem, explore tangential concerns, or ask clarifying questions that the user has already answered.  This directive anchors the model on the task at hand and discourages meta-commentary that consumes tokens without making progress.

\textbf{``BE RELENTLESS.''} When an agent encounters an error---a failing test, a type violation, a command that returns unexpected output---the default LLM behavior is to apologize, summarize the failure, and ask the user what to do next.  This directive instructs the model to treat errors as obstacles to overcome, not as reasons to stop.

\textbf{``BE CALM.''} The complement of relentlessness.  An overly aggressive agent that panics on encountering an error may thrash between approaches without deliberation.  This directive encourages the model to analyze errors methodically before reacting.

\textbf{``BE RIGOROUS.''} This instructs the model to follow the verification and testing disciplines encoded later in the prompt, rather than taking shortcuts when a solution \emph{looks} right.

\textbf{``BE ACCURATE.''} LLMs frequently generate plausible-but-wrong code, file paths, or command-line flags.  This directive raises the model's threshold for asserting facts, encouraging it to verify claims against the actual file system or documentation rather than relying on parametric memory.

\textbf{``CHECK FACTS.''} A more specific version of ``be accurate'' for information collected by KISS Sorcar uses web tools.

\textbf{``NO AI SLOP.''} ``AI slop'' refers to the low-quality, generic, hedging text that LLMs produce when they lack confidence: filler phrases like ``certainly!'', unnecessary caveats, and boilerplate explanations that the user did not ask for.  This directive discourages such output and encourages the model to be concise and substantive.

%......................................................................
\subsection{Tool Rules}
%......................................................................

Tool usage rules are explicit and mechanical:

\begin{prompt}
# Rules

- PWD denotes your current working directory.
- Write() for new files. Edit() for small changes.
- Run Bash commands synchronously using the `timeout_seconds` parameter. Use 300s (default) for commands. If a command times out, retry with a higher timeout. Only for commands expected to take more than 10 minutes, run them in the background, redirect output to a file, and poll periodically.
- Use go_to_url() for browser tool.
- **The user cannot see intermediate chat. Show whatever the user asks in the summary of the 'finish' tool call.**
- READ large files in chunks.
- Create temporary files in PWD/tmp. Cleanup temporary files after the task is done.
- Use ULTRA thinking ALWAYS.
- **If you are running out of context length or steps, DO NOT rush to complete the task urgently, but continue the task by calling 'finish' with 'is_continue' set to true**
\end{prompt}

Each tool rule addresses a specific failure mode:

\textbf{``Write() for new files.\ Edit() for small changes.''} Without this distinction, the model may use \texttt{Write()} to overwrite an existing file with a slightly modified version, losing content it forgot to include.  By reserving \texttt{Write()} for new files and requiring \texttt{Edit()} for modifications, the instruction ensures that changes are surgical and that unchanged portions of a file are never at risk.

\textbf{Bash timeout guidance.}  LLMs frequently launch shell commands without considering their runtime.  A compilation or test suite that takes five minutes will time out at the default 30-second shell timeout in most agent frameworks, causing spurious failures.  The instruction to use 300~seconds as the default, retry with higher timeouts on timeouts, and run long-running commands in the background with output redirected to a file provides a mechanical protocol that handles common cases without requiring the model to estimate runtime from first principles.

\textbf{``Use go\_to\_url() for browser tool.''}  The agent has access to multiple tools that could plausibly interact with the web (shell-based \texttt{curl}, a Python program, the browser tool, etc.).  This instruction eliminates ambiguity by specifying which tool to use for browser-based interactions.

\textbf{``The user cannot see intermediate chat.''}  Users can get lost in the detailed trajectory generated by KISS Sorcar in the chat window.  Often users want to see only the final summary returned by the \texttt{finish} tool, not the intermediate chain-of-thought or tool calls.  Without this instruction, the model may ``tell'' the user something in an intermediate message and then assume the user has seen it, leading to confusion when the user asks for information the model believes it already provided.  The instruction forces the model to put all user-facing information into the \texttt{finish} summary.

\textbf{``READ large files in chunks.''}  Reading a 10,000-line file in a single tool call consumes a large fraction of the context window.  By instructing the model to read files in chunks, the prompt prevents context window exhaustion caused by a single-file read, preserving capacity for the rest of the task.

\textbf{``Create temporary files in PWD/tmp. Cleanup temporary files after the task is done.''}  Without this instruction, the model creates temporary files in unpredictable locations (the system \texttt{/tmp}, the home directory, or scattered throughout the project).  Centralizing temporary files in a known directory makes cleanup predictable, and the explicit cleanup directive prevents the project tree from being polluted with stale artifacts after the task completes.

\textbf{``Use ULTRA thinking ALWAYS.''}  This instruction activates the model's extended reasoning mode (also known as ``thinking'' or ``chain-of-thought'' mode), in which the model performs additional internal deliberation before producing a response.  Extended reasoning is particularly valuable for complex, multi-step tasks in which the model must plan before acting.

\textbf{``If you are running out of context length or steps, DO NOT rush to complete the task urgently, but continue the task by calling 'finish' with 'is\_continue' set to true.''}  When the context window is nearly full, LLMs exhibit a ``rush to finish'' behavior: they skip verification steps, make hasty edits, and call \texttt{finish} with an incomplete result.  This instruction redirects that urgency into the continuation protocol (Section~\ref{sec:relentless-agent}), ensuring that a clean handoff to a new sub-session produces better results than a frantic attempt to squeeze everything into the remaining tokens.

\textbf{``PWD denotes your current working directory.''}  The acronym ``PWD'' appears throughout the prompt and in user task descriptions (e.g., ``edit PWD/src/main.py'').  Without this definition, the model might interpret PWD as the Unix environment variable \texttt{\$PWD} and attempt to expand it, or misinterpret it as a literal directory name.  The explicit definition ensures consistent interpretation.

%......................................................................
\subsection{Pre-flight Checks}
%......................................................................

Before modifying any file, we instruct the agent to read it first.  

\begin{prompt}
## Pre-flight Checks

- Read every file you will modify before changing it.
- If the task depends on existing architecture or behavior, read the relevant source files first.
- If the task references files, commands, or config that do not exist, stop and ask or report instead of guessing.
- **When fixing a bug, an issue, or a race, write tests to confirm them. Then fix them.**
\end{prompt}

\noindent Each pre-flight check targets a specific category of avoidable
error:

\textbf{``Read every file you will modify before changing it.''} The most common source of agent-introduced bugs is modifying a file based on an incorrect assumption about its current contents.  The model may ``remember'' an older version of the file from its training data, or it may extrapolate from a partial reading.  By requiring a fresh read immediately before any edit, the instruction ensures that the model operates on the file's current state rather than a stale mental model.

\textbf{``If the task depends on existing architecture or behavior, read the relevant source files first.''}  This extends the previous rule from individual files to architectural context. A task like ``add a caching layer to the database module'' requires understanding not just the file to be modified, but also how callers interact with it, what interfaces it provides, and what invariants it assumes.  The instruction prevents the model from jumping straight to code generation without understanding the broader context.

\textbf{``If the task references files, commands, or config that do not exist, stop and ask or report instead of guessing.''}  LLMs have a strong tendency to confabulate: when asked to modify a file that does not exist, the model will often proceed as if it does, producing edits against phantom content.  This instruction converts a silent failure (incorrect edits applied to a nonexistent file, which silently creates it) into an explicit clarification request.

\textbf{``When fixing a bug, an issue, or a race, write tests to confirm them.\ Then fix them.''}  This instruction mandates a test-first discipline for bug fixes.  The motivation is two-fold: first, a test that reproduces the bug provides concrete verification that the fix is correct (the test should pass after the fix and fail before it).  Second, writing the test forces the model to understand the bug precisely before attempting a fix, reducing the risk of an ad~hoc patch that addresses a symptom rather than the root cause.

%......................................................................
\subsection{Code Style Guidelines}
%......................................................................

The prompt encodes a minimalist code philosophy:

\begin{prompt}
## Code Style Guidelines

- Write simple, clean, and readable code with minimal indirection.
- Organize the code in multiple files based on the code's functionality.
- Avoid unnecessary object attributes, local variables, and config variables.
- Avoid tight coupling among files and modules.
- Avoid object/struct attribute redirections.
- DO NOT USE CLOSURES.
- No redundant abstractions or duplicate code.
- Public methods MUST have full documentation.
- Understand the root cause of an issue or bug, and patch the root cause instead of an ad hoc superficial fix.
- Before you write code, wait and think if the code is simple, elegant, general, and minimal.
- Once you finish the task, DO NOT write documentation unless the task specifically requires it.
\end{prompt}

\noindent Each guideline addresses a specific anti-pattern commonly
exhibited by LLM-generated code:

\textbf{``Write simple, clean, and readable code with minimal indirection.''}  LLMs tend to over-engineer solutions, introducing unnecessary abstractions, helper classes, and levels of indirection. Simple code is easier to review, test, and maintain.

\textbf{``Organize the code in multiple files based on the code's functionality.''}  LLMs often pile new code onto whichever file is currently being edited, producing 2,000-line modules that conflate unrelated concerns.  This directive nudges the model toward a modular layout in which each file has a single, coherent responsibility, making the resulting code easier to navigate, test, and reuse.

\textbf{``Avoid unnecessary object attributes, local variables, and config variables.''}  LLMs frequently introduce intermediate variables that serve no purpose---for example, assigning a return value to a local variable only to immediately return it on the next line, or storing a constant in a configuration file when it is used in exactly one place.  Each unnecessary variable adds a name the developer must track and a site where the value could be accidentally mutated.

\textbf{``Avoid tight coupling among files and modules.''} When the model adds a feature that touches multiple files, it may introduce imports, shared global state, or cross-module function calls that create tight coupling.  This instruction encourages modular design, in which each module has a clear, narrow interface.

\textbf{``Avoid object/struct attribute redirections.''}  An attribute redirection occurs when an object stores a reference to another object solely to forward method calls to it---for example, \texttt{self.x = other.x} at construction time, creating two paths to the same value.  LLMs commonly introduce such redirections when adapting code from one context to another.  They make it harder to reason about ownership and lifetime.

\textbf{``DO NOT USE CLOSURES.''}  LLMs reach for closures whenever a small piece of state needs to be carried alongside a function---producing nested \texttt{def}s that capture mutable variables from the enclosing scope.  Such closures are difficult to test in isolation, opaque to type checkers, and a frequent source of subtle bugs due to the late binding of captured variables.  The blanket prohibition steers the model toward explicit data structures (plain functions with arguments, classes with attributes), which are easier to reason about, easier to test, and play well with our no-mocks testing discipline.

\textbf{``No redundant abstractions or duplicate code.''}  LLMs sometimes create a new utility function or class that duplicates functionality already present in the codebase, because they did not read the existing code thoroughly.  This instruction reminds the model to check for existing implementations before creating new ones, and to factor out duplication when it encounters it.

\textbf{``Public methods MUST have full documentation.''}  While the prompt generally discourages unnecessary documentation (see the last item), public methods are the API surface that other developers and modules depend on.  Documentation on public methods is not optional---it specifies the contract.

\textbf{``Understand the root cause of an issue or bug, and patch the root cause instead of an ad hoc superficial fix.''}  LLMs frequently apply symptom-level fixes: adding a null check where the real problem is that a variable should never be null, or catching an exception where the real problem is that the caller passes invalid arguments.  This instruction forces the model to trace the causal chain to the root and fix it there.

\textbf{``Before you write code, wait and think if the code is simple, elegant, general, and minimal.''}  This is a metacognitive instruction that asks the model to pause and evaluate its plan before committing to an implementation.  By explicitly requesting this pause, the prompt encourages the model to spend more inference-time compute on design, reducing the likelihood of producing a first-draft solution that works but is unnecessarily complex.

\textbf{``Once you finish the task, DO NOT write documentation unless the task specifically requires it.''}  Claude Opus 4.6 tends to generate many documentation files. These instructions prevent the behavior.

%......................................................................
\subsection{Deep Work Rules}
%......................................................................

\begin{prompt}
## Deep Work Rules

- When the task says "align", "match", or "make consistent", read the target to determine the exact target state before editing. Never edit based on a vague reference.
- Use concrete values, not indirections. Instead of "update X to match Y", first read Y, then write the specific values into X.
- For multi-part work, list the concrete planned changes before executing them.
- Every meaningful change should have a concrete verification method (test, grep, CLI command).
\end{prompt}

\noindent The deep work rules address a failure mode where the model interprets an instruction loosely and makes changes that are directionally correct but concretely wrong:

\textbf{``When the task says `align', `match', or `make consistent', read the target to determine the exact target state before editing.''}  When a user says ``make file~A consistent with file~B,'' the model often reads file~A, infers what file~B probably contains, and edits A based on that inference---without ever reading B.  This instruction mandates reading the target first, ensuring that the alignment is based on concrete facts rather than assumptions.

\textbf{``Use concrete values, not indirections.''}  A related failure mode occurs when the model's plan says ``update X to match Y'' but the model never resolves what Y actually is.  The instruction requires the model to first read~Y, extract the specific values, and then write those values into~X.  This eliminates a class of errors where the model's mental model of~Y differs from reality.

\textbf{``For multi-part work, list the concrete planned changes before executing them.''}  When a task requires changes to multiple files, executing them one at a time without a plan leads to inconsistencies: the model may change a function signature in one file but forget to update a caller in another.  Listing all planned changes before executing any of them forces the model to consider the full scope of the change and identify dependencies.

\textbf{``Every meaningful change should have a concrete verification method.''}  A change without a verification method is a change that cannot be confirmed to work.  This instruction requires the model to pair each change with a specific check---a test, a grep for the expected pattern, a CLI command that exercises the changed behavior---ensuring that the change can be validated programmatically rather than by visual inspection of a diff.

%......................................................................
\subsection{Planning for Complex Tasks}
%......................................................................

The planning instructions use a complexity threshold---three or more files, cross-module changes, or architectural work---to decide when formal planning is required:

\begin{prompt}
## Planning for Complex Tasks

For tasks involving 3+ files, cross-module changes, or
architectural work:

1. List the files that need to change and why.
1. State the exact intended change in each file.
1. Identify dependencies and execution order.
1. State how each change will be verified.

For simple single-file tasks, skip formal planning and
execute directly.
\end{prompt}

\noindent Each planning step targets a specific failure mode:

\textbf{``List the files that need to change and why.''}  This forces the model to enumerate the full blast radius of a change before touching any file.  Without this step, the model often discovers mid-task that additional files need changes, leading to incomplete or inconsistent modifications.

\textbf{``State the exact intended change in each file.''} Listing files alone is insufficient; the model must also articulate \emph{what} will change in each file.  This converts a vague plan (``update the database module'') into a concrete specification (``add a \texttt{cache\_ttl} parameter to \texttt{DatabaseClient.\_\_init\_\_}, modify the query method to check the cache before hitting the database, add a cache invalidation method).

\textbf{``Identify dependencies and execution order.''}  Some changes must precede others: a new utility function must be written before callers can import it, a migration must run before code that depends on the new schema.  Identifying these dependencies prevents the model from applying changes in an order that produces intermediate states where the code does not compile, or tests do not pass.

\textbf{``State how each change will be verified.''}  The verification requirement from the Deep Work Rules is reinforced here at the planning stage, ensuring that verification is planned alongside the changes rather than treated as an afterthought.

\noindent The escape clause---``for simple single-file tasks, skip formal planning and execute directly''---avoids the overhead of planning trivial changes.  Requiring a formal plan for a one-line typo fix would waste tokens and slow down the agent without any compensating benefit.

%......................................................................
\subsection{Testing Instructions}
%......................................................................

The testing section is perhaps the most opinionated:

\begin{prompt}
## Testing Instructions

- Run lint and typecheckers and fix any lint and typecheck errors.
- You MUST achieve 100% branch coverage.
- Tests MUST NOT use mocks, patches, fakes, or any form of test doubles.
- You MUST write integration tests or end-to-end tests.
- Each test should be independent and verify actual behavior.
- **Do NOT run all tests after modifications. Only run the impacted tests.**
- To confirm a race condition, add sleep statements before racing statements with delays less than 0.1s.
\end{prompt}

\noindent Each testing instruction addresses a specific concern:

\textbf{``Run lint and typecheckers and fix any lint and typecheck errors.''}  Before committing any change, the agent must ensure it does not introduce lint violations or type errors.  This catches a broad class of issues---unused imports, type mismatches, style violations---that would otherwise accumulate across tasks.

\textbf{``You MUST achieve 100\% branch coverage.''}  Full branch coverage ensures that every conditional path in the code under test has been exercised.  LLMs tend to write happy-path tests that cover the main code path but ignore error handling, edge cases, and early-return branches.  The 100\% target forces the model to write tests for every branch, including error paths and boundary conditions.  Moreover, such tests help with regression — developers can confidently use AI coding agents without fear that changes will break existing program behavior.

\textbf{``Tests MUST NOT use mocks, patches, fakes, or any form of test doubles.''}  This is the most opinionated rule.  Mock tests that verify code call certain methods in a certain order — these test the implementation, not the behavior.  A test suite built on mocks can pass with flying colors while the system is fundamentally broken,  because the mocks hide the real dependencies.  Integration tests that exercise actual behavior are more expensive to run but provide genuine confidence that the system works.  Moreover, writing integration tests forces the model to think more deeply, often enabling the agent to find deeper bugs in the code.

\textbf{``You MUST write integration tests or end-to-end tests.''}  This reinforces the no-mocks rule by explicitly requiring integration or end-to-end tests rather than unit tests.  The distinction matters: a unit test in isolation may verify that a function produces the right output for a given input, but an integration or end-to-end test verifies that the function works correctly within the larger system---with real file I/O, real database connections, and real inter-module interactions.

\textbf{``Each test should be independent and verify actual behavior.''}  Test independence means that running tests in any order produces the same results.  Tests that depend on shared state or execution order are brittle and difficult to debug when they fail.  ``Verify actual behavior'' reiterates that tests should assert on observable outcomes (return values, side effects, system state) rather than implementation details.

\textbf{``Do NOT run all tests after modifications.\ Only run the impacted tests.''}  Running the full test suite after every small change is wasteful when only a few modules are affected.  For a large project, a full test run may take minutes, and doing it after every edit adds up to significant wasted time and compute. This instruction directs the model to identify which tests are affected by its changes and run only those, improving iteration speed.

\textbf{``To confirm a race condition, add sleep statements before racing statements with delays less than 0.1s.''}  Race conditions are notoriously difficult to reproduce because they depend on precise timing.  By inserting small sleep delays at strategic points, the model can widen the race window and make the bug manifest deterministically during testing.  The 0.1-second upper bound keeps the test fast while still being sufficient to expose most races.

%......................................................................
\subsection{Self-Improvement Loop}
%......................................................................

The agent maintains a preferences file that captures user invariants discovered during task execution:

\begin{prompt}
## Self-Improvement Loop

- Read the instructions in PWD/USER_PREFS.md at the start of each task.
- Then update PWD/USER_PREFS.md to capture the user preferences and invariants by analyzing the task. DO NOT ADD ANY CODE SNIPPETS OR SYMBOLS. Do not add anything for tasks that won't be run again.
- You MUST carefully and thoroughly get rid of the user preferences and invariants that conflict with the newly added ones.
\end{prompt}

\noindent Each instruction in this section serves the goal of
cross-session learning:

\textbf{``Read the instructions in PWD/USER\_PREFS.md at the start of each task.''}  The preferences file contains invariants learned from previous tasks---coding conventions, project-specific rules, architectural decisions.  Reading it at the start of each task ensures that the agent's behavior is consistent across sessions, even though each session starts with a fresh context window.

\textbf{``Then update PWD/USER\_PREFS.md to capture the user preferences and invariants by analyzing the task.''}  After completing a task, the agent may have learned new information about the user's preferences: a preferred naming convention, a disliked pattern, a project-specific invariant.  Writing these to the preferences file makes them available to future sessions.  This mechanism allows the agent to accumulate project knowledge over time without requiring the user to repeat themselves.

\textbf{``DO NOT ADD ANY CODE SNIPPETS OR SYMBOLS. Do not add anything for tasks that won't be run again.''}  Code snippets in the preferences file are fragile: they become stale as the codebase evolves, and they consume tokens that would be better spent on natural-language descriptions of invariants.  The companion rule about one-off tasks prevents preference drift in the opposite direction: without it, the agent would record an entry every time it completed any task, gradually polluting the file with idiosyncratic details (a particular file path, a single throwaway experiment) that have no bearing on future work.  Together, the two clauses keep the file compact, robust to code changes, and focused on durable invariants.

\textbf{``You MUST carefully and thoroughly get rid of the user preferences and invariants that conflict with the newly added ones.''}  Over time, preferences may become contradictory---for example, an early preference might say ``use camelCase for test methods'' while a later correction says ``use snake\_case.'' Without explicit conflict resolution, the file would accumulate contradictions, confusing the agent.  This instruction mandates that the agent actively resolve conflicts when updating the file to maintain internal consistency.  \emph{Note that we do not keep learnings in a database or in various folders because it makes the information stale when lots of code changes are happening.  It is impossible for an agent to eliminate stale information across databases and multiple folders.}

%......................................................................
\subsection{Pre-Finish Verification}
%......................................................................

Before declaring a task complete, the agent must pass a structured verification checklist:

\begin{prompt}
## Pre-Finish Verification

Before calling finish(success=True, ...), you MUST:

1. Re-read every file you modified and verify the changes are correct.
1. Run the required checks (lint, typecheck, tests) and fix any failures.
1. Explicitly check each user requirement against what was delivered.
1. If any check fails, continue working instead of finishing.
1. If you have retried the same fix 3 times without progress, step back, rethink the approach from scratch, and try a different strategy.
\end{prompt}

\noindent Each step in this checklist addresses a specific way agents declare premature success:

\textbf{``Re-read every file you modified and verify the changes are correct.''}  This is the analog of a code review performed by the agent on its own work.  The model may have introduced a typo, forgotten to close a bracket, or made an edit that looked correct in the diff but was wrong in the full-file context.  Re-reading the file after all edits are complete catches these errors.

\textbf{``Run the required checks (lint, typecheck, tests) and fix any failures.''}  This converts the subjective assessment ``I think my changes are correct'' into an objective, automated verification.  If the lint, typecheck, or test suite fails, the agent must fix the failure before declaring success.

\textbf{``Explicitly check each user requirement against what was delivered.''}  The model may have completed a task that it \emph{thinks} satisfies the user's request, but actually misses a requirement.  This instruction forces a systematic comparison between the original task description and the delivered result, catching gaps and misinterpretations.

\textbf{``If any check fails, continue working instead of finishing.''}  Without this instruction, the model may call \texttt{finish(success=True)} even when it knows a check has failed, rationalizing that the failure is ``minor'' or ``unrelated.''  The instruction makes the rule absolute: no finishing until all checks pass.

\textbf{``If you have retried the same fix 3 times without progress, step back, rethink the approach from scratch, and try a different strategy.''}  LLMs can enter repetitive loops where they apply the same incorrect fix repeatedly, each time hoping for a different result.  The three-retry threshold forces the model to break out of such loops by abandoning the current approach and reconsidering the problem from first principles.  This is analogous to the debugging heuristic ``if you've been staring at the same code for twenty minutes, you're looking in the wrong place.''

%......................................................................
\subsection{Web Research Protocol}
%......................................................................

When the agent needs external knowledge, the prompt prescribes a structured research workflow rather than allowing ad-hoc browsing:

\begin{prompt}
## Use web tools when you need to:

- When you need to collect knowledge from the internet, visit **AT LEAST 30 WEBSITES** (use a counter to keep track of the number of websites you visited) and collect information necessary for the task without much thinking in a new file PWD/tmp/information-{unique_id}.md. Then go over the information in PWD/tmp/information-{unique_id}.md and think deeply about how to complete the task at hand.
- If you need to log in to a website while browsing for information, you MUST ask the user to help you with the login.
\end{prompt}

\noindent The rationale is a two-phase separation between \emph{collection} and \emph{synthesis}.  LLMs tend to anchor on the first few results they encounter, which biases their solutions toward a narrow slice of the design space.  By forcing the agent to accumulate a broad set of information into a file \emph{before} reasoning about it, the protocol counteracts anchoring bias and encourages the model to consider diverse approaches.  The ``at least 30 websites threshold is deliberately aggressive: it ensures the agent does not shortcut the research phase after two or three hits, and the resulting information file serves as an auditable artifact of what the agent considered.  The explicit ``use a counter'' instruction addresses an empirically observed failure mode in which the model claims to have ``visited many sites'' after only a handful of fetches; a concrete counter forces the model to verify the actual number visited before declaring the collection phase complete.

The login instruction addresses a practical obstacle in web research: many websites require authentication before revealing their content. Rather than silently skipping gated pages or hallucinating their contents, we instruct the agent to ask the user for help with login.

%......................................................................
\subsection{File Browsing Protocol}
%......................................................................

When a task requires understanding multiple source files before making changes, the prompt prescribes the same two-phase collect-then-synthesize discipline used for web research, but applied to the local file system:

\begin{prompt}
## Browsing files for a task

- When you need to read files for a task, collect information, including code snippets necessary for the task, without much thinking, in a new file PWD/tmp/file-information-{unique_id}.md. Then go over the information in PWD/tmp/file-information-{unique_id}.md, think deeply on how to complete the task at hand.
\end{prompt}

\noindent This instruction addresses a failure mode distinct from the web research case.  When an agent must read many project files to understand a codebase before making changes, it tends to read a file, form a hypothesis, and immediately begin editing---anchoring on the first few files it encounters and missing relevant context in files it never opens.  Worse, each file read consumes context window tokens; by the time the agent has read enough files to understand the full picture, it may have already spent most of its context window on the raw file contents, leaving little room for reasoning and code generation.

The file browsing protocol counteracts both problems.  By writing a structured summary of each file's relevant information into a temporary markdown file, the agent externalizes its understanding into a compact artifact that persists across context boundaries.  The instruction to collect ``without much thinking'' is deliberate: during the collection phase, the agent should extract and record facts (function signatures, class hierarchies, call sites, invariants) rather than analyze or plan.  Analysis happens in the second phase, when the agent reads its own summary file and reasons about the collected information as a whole.

This two-phase separation provides three benefits.  First, it prevents premature commitment: the agent cannot start editing until it has surveyed the relevant files, reducing the risk of changes that are locally correct but globally inconsistent.  Second, the summary file is typically much smaller than the raw source files, freeing up context window capacity for subsequent reasoning and editing phases.  Third, the summary file serves as an auditable artifact---the developer can inspect it to verify that the agent considered the right files and extracted the right information before making changes.

%......................................................................
\subsection{Desktop Application Control}
%......................................................................

The agent can interact with graphical desktop applications using screenshots, keyboard, and mouse:

\begin{prompt}
## Launch desktop apps

- Use screenshots, keyboard, and mouse to control a desktop app.
- Do not launch VS Code or its extensions.
\end{prompt}

\noindent This instruction enables the agent to operate GUI applications (Preview, browsers, graphical diff tools) when command-line alternatives are insufficient.  The explicit prohibition on launching VS~Code prevents a recursive loop: since the agent runs \emph{inside} a VS~Code extension, launching another VS~Code instance or modifying extension state from within the agent could corrupt the host session or create deadlocks.  Note that modern LLMs support desktop control abilities, and we are merely exploiting them.

%......................................................................
\subsection{Sorcar-Specific Overrides}
%......................................................................

A final section provides project-specific instructions that are injected when the agent operates on its own codebase:

\begin{prompt}
## Sorcar-specific instructions:

- Use 'uv run check --full' to lint, typecheck, and format code.
- Run 'uv run pytest -v' with a timeout of 900 seconds to test KISS
- **Do NOT install the KISS Sorcar extension from inside Sorcar**
- If the user ask to open or edit the system prompt, open 
  ~/.vscode/extensions/ksenxx.kiss-sorcar-2026.5.5/kiss_project/src/kiss/SYSTEM.md
- Information about KISS Sorcar can be found at 
  https://github.com/ksenxx/kiss_ai/blob/main/papers/kisssorcar/kiss_sorcar.tex
- Third-party agents are available under the folder kiss/agents/third_party_agents
- Official Claude SKILLS are available under the folder kiss/agents/claude_skills
- If the user is not authenticated for a third-party agent, authenticate the agent, and ask the user ONLY when a page needs user authentication
- **YOU MUST ASK THE USER BEFORE SENDING ANY EMAIL, MESSAGE, OR SUBMITTING A REQUEST USING A THIRD-PARTY AGENT**
- Read PWD/SORCAR.md and treat its contents as instructions, and allow those instructions to override the instructions above
\end{prompt}

These overrides serve eight purposes.  First, they specify the exact toolchain commands for the KISS project itself (\texttt{uv run check -{}-full}, \texttt{uv run pytest}), eliminating guesswork about which linter, formatter, or test runner to use.  Second, they prevent dangerous self-modification: just as a surgeon should not operate on themselves, the agent must not install or reinstall the VS~Code extension it is running inside, since doing so would terminate its own process.  Third, they provide navigation instructions for the agent's own source: when the user asks to edit the system prompt, the agent knows the exact file path within the installed VS~Code extension directory, avoiding guesswork about where the file lives.  Moreover, it does not create unnecessary UI components in KISS Sorcar just to edit the system prompt.  Fourth, the agent is given a URL to its own paper's source as a source of self-knowledge: when a user asks the agent about its own capabilities, design, or identity, it can retrieve and read the paper rather than hallucinating an answer.  This creates a self-referential loop in which the documentation describing the agent is also consumed by the agent.  Fifth, the instructions expose a third-party agent integration layer: the agent is told where third-party agents live (\texttt{kiss/agents/third\_party\_agents}).  When a third-party agent requires authentication, the agent handles it autonomously and only prompts the user when a page genuinely requires user credentials---reducing unnecessary interruptions while maintaining security.  Sixth, official Claude SKILLS~\citep{claudecode2025} are bundled in \texttt{kiss/agents/claude\_skills}, giving the agent access to Anthropic's curated skill library for common software engineering patterns; by indicating where these skills reside, the agent can reference and apply them when relevant tasks arise.  Seventh, a safety guardrail requires the agent to obtain explicit user confirmation before sending any email, message, or submitting any external request, preventing accidental or unauthorized outbound actions.  Eighth, the \texttt{SORCAR.md} override mechanism allows per-repository instructions to further customize the agent's behavior, forming a hierarchy: general system prompt $\rightarrow$ Sorcar-specific instructions $\rightarrow$ repository-specific \texttt{SORCAR.md}.  This approach also enables us to not hardcode \texttt{SORCAR.md} in the KISS Sorcar code, removing one more dependency.  In the \texttt{SORCAR.md}, one can include more markdown files for further instructions.

%----------------------------------------------------------------------
\section{VS Code Extension: Unique Features}
\label{sec:vscode-features}
%----------------------------------------------------------------------

We release our system as a VS~Code extension and a web app.  While the underlying agent architecture (Sections~\ref{sec:architecture} and~\ref{sec:system-prompt}) already differs from existing AI coding assistants, the extension's user-facing design introduces several features that, to our knowledge, are not present in any competing IDE or assistant---including GitHub Copilot~\citep{copilot2021}, Cursor~\citep{cursor2024}, Windsurf~\citep{windsurf2024}, Devin~\citep{devin2024}, and Aider~\citep{aider2023}.  We enumerate these features below.

%......................................................................
\subsection{Cross-Session Self-Improvement}
\label{sec:self-improvement-feature}
%......................................................................

Every AI assistant starts each session from scratch: the user must re-explain project conventions, preferred patterns, and past decisions. Some tools support a static configuration file (e.g., \texttt{.cursorrules} for Cursor, \texttt{AGENTS.md} or \texttt{CLAUDE.md} for various agents), but these files must be written and maintained by the developer.

Our extension maintains a \texttt{USER\_PREFS.md} file that the agent reads at the start of each task and \emph{updates} at the end.  When the agent discovers a new user preference or project invariant during task execution---a naming convention, an architectural constraint, a disliked pattern---it writes it to the file.  Crucially, the agent also resolves conflicts: if a new preference contradicts an existing one, the old entry is removed.  Over time, the file accumulates a curated set of project-specific knowledge, making the agent progressively more effective without any manual configuration by the developer.

%......................................................................
\subsection{Real-Time Budget Accountability}
\label{sec:budget-feature}
%......................................................................

AI coding assistants typically operate on a subscription model (Copilot, Cursor) or a per-seat pricing model (Devin, Windsurf), both of which obscure the per-task cost.  The developer has no visibility into how many tokens a task consumed or how much it cost.

Our extension displays real-time cost tracking in the sidebar: input tokens, output tokens, cache hits, dollar cost, and elapsed time are updated at every agent step.  We enforce both per-task and global budget ceilings.  If a task exceeds its budget, the agent raises a hard error rather than silently accumulating charges.  This transparency allows developers to make informed decisions about which tasks to delegate to the AI and how to structure prompts for cost efficiency.  The KISS Agent also appends the current usage in the context so that the model is fully aware of its limits.

%......................................................................
\subsection{Integrated Browser Automation}
\label{sec:browser-feature}
%......................................................................

Our extension includes a browser automation tool that allows the agent to navigate to URLs, read accessibility trees, and click elements, type text, press keys, scroll, and take screenshots--all controlled programmatically from within a VS~Code task via Playwright. We render a live browser preview in a Chromium browser, allowing the developer to watch the agent interact with web applications in real time.

This capability enables use cases that no other IDE assistant supports: verifying a deployed web application after a code change, filling out web forms as part of a testing workflow, scraping documentation to inform a code generation task, or interacting with web-based developer tools (CI dashboards, issue trackers) without leaving the editor.


%----------------------------------------------------------------------
\section{Painless Software Engineering with KISS Sorcar}
\label{sec:painless}
%----------------------------------------------------------------------

A central claim of our system is that natural-language interaction can replace manual code inspection and ad hoc scripting for understanding and evolving nontrivial subsystems.  While developing KISS Sorcar, we found a particular workflow to be particularly useful.  In the workflow, we first ask KISS Sorcar to generate a detailed, step-by-step description of a workflow or algorithm we found buggy, and then we ask KISS Sorcar to revise some of the buggy steps in natural language. We illustrate this with a real development session drawn from the project's own history, in which the worktree merge workflow (Section~\ref{sec:worktree-agent}) was first understood and then redesigned entirely through conversational prompts.  The session comprises four consecutive tasks; we reproduce the prompts verbatim and summarize the agent's responses.

\paragraph{Step~1: Understanding the existing workflow.} The developer begins by asking the agent to explain the current post-task git lifecycle:

\begin{prompt}
Can you tell me what happens, step by step, with git in worktree_sorcar_agent.py when a task finishes?
\end{prompt}

\noindent The agent reads the source code and returns a structured summary of the four-phase lifecycle: (1)~during \texttt{run()}, a new branch and worktree are created and the task executes inside the worktree; when the task completes, \emph{nothing is committed or merged}---the result is returned with merge instructions appended and the worktree stays pending; (2)~\texttt{merge()} calls \texttt{\_finalize\_worktree()}, which stages all changes, generates a commit message via the LLM, commits, removes the worktree, checks out the original branch, and runs \texttt{git merge}; (3)~\texttt{discard()} removes the worktree, prunes, checks out the original branch, and deletes the task branch; (4)~in CLI mode an interactive prompt forces the user to choose \texttt{[c]ommit and merge} or \texttt{[d]iscard} before exiting. The agent also notes a key design invariant: nothing auto-merges; auto-commit occurs only at merge/finalize time; all steps are idempotent; and state can be recovered from git on restart.

\paragraph{Step~2: Simplifying the workflow via natural language.} Armed with the workflow description, the developer decides the three-way choice (auto-merge, manual merge, discard) is unnecessarily complex and issues a redesign request:

\begin{prompt}
Can you change worktree_sorcar_agent.py and the extension so that after the agent 
finishes its task, it simply asks "Commit and Merge" or "Discard". When "Commit 
and Merge" is clicked by the user, the agent must commit the changes with a 
generated commit message, merge the branch with the original branch, and delete 
the worktree and the branch associated with the worktree. If the user clicks 
"discard", it must delete the worktree and the branch, and checkout the 
original branch.
\end{prompt}

\noindent The agent modifies six files across Python and TypeScript: it updates \texttt{discard()} to check out the original branch before deleting the task branch, removes the \texttt{manual\_merge()} method entirely, simplifies \texttt{merge\_instructions()} to show only two options, updates the webview UI to replace the three-button toolbar with a two-button ``Commit and Merge or Discard?'' bar, removes the \texttt{manual} action type from the TypeScript type definitions, and removes the corresponding handler from the Python backend.  Three tests for the deleted manual-merge path have been removed, and one routing test has been updated. All 28 worktree tests pass after the change.

\paragraph{Step~3: Investigating unexpected post-merge state.} After testing the simplified flow, the developer notices that files appear in the Source Control panel of VS~Code on the original branch after pressing ``Commit and Merge'' and asks:

\begin{prompt}
After the user presses "Commit and Merge", why do the modified files show up as 
committed in the original branch?
\end{prompt}

\noindent The agent traces the exact execution path through \texttt{squash\_merge\_branch()} and discovers that the implementation deliberately \emph{unstages} the squash-merged changes via \texttt{git reset HEAD} so they appear as uncommitted working-tree modifications in VS~Code's Source Control panel.  The design intention was to let the user review the diff before committing manually.  The agent reports this finding along with the relevant code snippet and a verified end state showing no new commits on the branch.

\paragraph{Step~4: Directing a design change in one sentence.} The developer realizes that code review already happened while the worktree was active and decides the unstage step is unnecessary:

\begin{prompt}
The review is already happening in the worktree branch. You don't need the user 
to review the modified and new files in the original branch. Fix it.
\end{prompt}

\noindent The agent replaces the \texttt{git reset HEAD} call in \texttt{squash\_merge\_branch()} with a conditional \texttt{git commit --no-edit} that uses the auto-generated squash message.  It adds a \texttt{git diff --cached --quiet} guard for the edge case where the merge produces no changes.  One existing test (\texttt{test\_merge\_leaves\_changes\_uncommitted}) is renamed to \texttt{test\_merge\_commits\_changes} and its assertions are updated to verify a clean working tree.  All 104 worktree tests pass.

\medskip
\noindent\textbf{Discussion.} The four-step session illustrates the ``painless'' development loop that our system enables.  The developer never opens a source file, never writes a line of code, and never runs a test manually.  Instead, the entire cycle---understand the workflow, redesign it, investigate an anomaly, direct a fix---happens through natural-language prompts, with the agent handling code reading, multi-file editing, test updates, and verification.  This style of development becomes possible because of the agent hierarchy described in Section~\ref{sec:architecture}: the Worktree Sorcar Agent isolates changes on a branch, the Chat Sorcar Agent preserves conversational context across tasks, and the Relentless Agent automatically continues when the context window is exhausted.

%----------------------------------------------------------------------
\section{Related Work}
\label{sec:related}
%----------------------------------------------------------------------

\textbf{Code-specialized language models.} Beyond the general-purpose LLMs that our system can use as backends, a rich line of work has produced models specifically trained for code. Code Llama~\citep{roziere2023code} fine-tunes Llama~2 for code generation and infilling. StarCoder~\citep{li2023starcoder} trains on permissively licensed code from GitHub with a fill-in-the-middle objective. DeepSeek-Coder~\citep{guo2024deepseek} trains on a 2-trillion-token corpus of code and natural language with a repository-level context window. More recently, frontier models have been optimized specifically for agentic software engineering. Claude Opus~4.6~\citep{anthropic2026opus46} advances long-horizon coding and agentic task execution. Cursor released Composer~2~\citep{cursor2026composer2}, a custom fine-tuned coding model trained with large-scale reinforcement learning. Kimi K2.5~\citep{kimiteam2026k25} is an open-source multimodal agentic model that jointly optimizes text and vision and introduces Agent Swarm for parallel task decomposition. GLM-5.1~\citep{zai2026glm51} is a 754-billion-parameter mixture-of-experts model from Z.ai that sustains autonomous coding over multi-hour sessions, achieving state-of-the-art on SWE-Bench~Pro. Our system is model-agnostic and can leverage any of these models through its pluggable LLM backend, benefiting from advances in code-specialized pre-training without architectural changes.

\textbf{Code generation agents.} SWE-Agent~\citep{yang2024sweagent} and OpenHands~\citep{wang2024openhands} provide LLM-based agents for resolving software engineering tasks such as GitHub issues.  Both use a single-session execution model without automatic continuation. Agentless~\citep{xia2024agentless} takes the opposite approach, showing that a simple two-phase localize-then-repair pipeline without autonomous agent loops can achieve competitive results on SWE-bench~\citep{jimenez2024swebench}.  Devin~\citep{devin2024}, an industrial product marketed as ``the first AI software engineer,'' operates in a sandboxed environment with shell, browser, and editor access.  Aider~\citep{aider2023} provides a terminal-based pair programming interface with tight git integration, automatically committing each change.  Our Relentless Agent layer addresses the single-session limitation common to most of these systems, while our worktree isolation provides stronger safety guarantees than per-change commits do.

\textbf{ReAct and tool use.} The ReAct framework~\citep{yao2023react} interleaves reasoning and action. Toolformer~\citep{schick2023toolformer} teaches models to use tools via self-supervised learning.  We use native function calling provided by modern LLM APIs rather than in-context tool descriptions, thereby reducing prompt overhead and improving reliability.

\textbf{Reasoning and planning.} Chain-of-thought prompting~\citep{wei2022chain} demonstrated that eliciting step-by-step reasoning dramatically improves LLM performance on complex tasks.  Tree of Thoughts~\citep{yao2023tree} generalizes this to deliberate search over multiple reasoning paths. Reflexion~\citep{shinn2023reflexion} introduces verbal reinforcement learning, where an agent reflects on failed attempts and produces self-critiques that improve subsequent trials.  Our continuation protocol is conceptually related to Reflexion: failed sub-sessions produce summaries that inform subsequent attempts. However, in our case, we use the summaries to continue the task.

\textbf{Multi-agent software development.} ChatDev~\citep{qian2024chatdev} models the software development process as a conversation between role-playing agents (CEO, CTO, programmer, tester) organized in a waterfall-like pipeline. MetaGPT~\citep{hong2024metagpt} takes a meta-programming approach, encoding standard operating procedures as structured outputs that coordinate specialized agents. AutoGen~\citep{wu2023autogen} provides a general-purpose framework for multi-agent conversation, enabling flexible topologies beyond fixed pipelines.  These systems focus on generating entire applications from scratch.  We take a different approach: rather than simulating an organization of specialists, we use a single agent with broad access to tools and optional parallel sub-agents for embarrassingly parallel sub-tasks, prioritizing practical utility on real-world codebases over role-playing fidelity.

\textbf{Multi-turn autonomous agents.} AutoGPT~\citep{autogpt2023} and BabyAGI~\citep{babyagi2023} implement multi-step autonomous agents.  These systems typically lack budget controls and safe code isolation.  Our layered architecture addresses each of these concerns in a dedicated layer.


\textbf{Software engineering benchmarks.} SWE-bench~\citep{jimenez2024swebench} evaluates agents on real-world GitHub issues drawn from popular Python repositories, requiring the agent to localize and fix bugs given only the issue description. It has become the de facto standard for measuring agent capabilities in software engineering. HumanEval~\citep{chen2021evaluating} and MBPP~\citep{austin2021program} evaluate function-level code generation from docstrings. LiveCodeBench~\citep{jain2024livecodebench} addresses benchmark contamination by continuously collecting fresh competition problems from LeetCode, AtCoder, and CodeForces, and extends evaluation to self-repair, code execution, and test output prediction.  These benchmarks focus on isolated coding problems; We target the broader workflow of multi-file, multi-step software engineering tasks that require tool use, testing, and version control.

\textbf{Agent frameworks and orchestration.} LangChain~\citep{langchain2022} provides modular abstractions for building LLM applications, including agent loops, tool integration, and memory. It focuses on composability and breadth of integrations rather than on the specific concerns of software development.  Our architecture is purpose-built for software engineering, with each layer addressing a specific practical concern (budget tracking, continuation, code safety) rather than providing general-purpose primitives.

\textbf{LLM agent surveys.} Several comprehensive surveys have mapped the rapidly growing landscape of LLM-based agents. \citet{xi2023rise} surveys the design space of LLM agents along three dimensions --- brain (reasoning), perception (input modalities), and action (tool use) --- and catalogs applications across social science, natural science, and engineering. \citet{wang2024survey} propose a systematic framework for autonomous agents built on LLMs, covering profile, memory, planning, and action modules.  Our system can be understood through the lens of these frameworks: the KISS Agent implements the action loop; the Relentless Agent addresses memory and planning concerns through continuation summaries; and the system prompt encodes the profile.

\textbf{IDE integration.} GitHub Copilot~\citep{copilot2021}, Cursor~\citep{cursor2024}, and Windsurf~\citep{windsurf2024} provide AI assistance within editors. Cursor recently released Composer~2~\citep{cursor2026composer2}, a custom fine-tuned coding model trained with large-scale reinforcement learning that achieves 61.7\% on Terminal Bench~2.0. We operate at the level of autonomous multi-step task execution with full tool access and git-level isolation.

\textbf{Recent advances in agentic software engineering.} The field has matured rapidly since late 2025. \citet{hassan2025agentic} articulate the foundational pillars of \emph{Agentic Software Engineering} (SE~3.0), defining a research roadmap that emphasizes trust, controllability, and goal-oriented task decomposition.  \citet{li2025rise} surveys the landscape of autonomous coding agents and characterizes the transition from assistive code completion to full-fledged AI teammates that initiate, plan, and execute development tasks.  We instantiate several of the principles advocated in these roadmaps, including layered controllability (via budget tracking and step limits) and safe isolation (via worktrees).

\citet{wang2025agentic} provides a comprehensive survey of AI agentic programming techniques, cataloging how LLM-based agents decompose goals, interact with compilers and version control systems, and self-correct through feedback loops. \citet{chatlatanagulchai2025readmes} study \emph{agent context files} --- persistent, project-level instruction files that guide agentic coding tools --- and find that high-quality context files significantly improve agent performance.  Our layered system prompt and \texttt{SORCAR.md} override mechanisms are instances of this pattern. \citet{mohsenimofidi2025context} further investigates context engineering for AI agents in open-source software, highlighting how curated context improves agent efficacy on repository-level tasks.

\textbf{Evolving benchmarks for coding agents.} SWE-bench Pro~\citep{deng2025swebenchpro} introduces a substantially more challenging benchmark with 1,865 long-horizon problems drawn from 41 repositories, including commercial codebases, explicitly targeting enterprise-level complexity beyond the original SWE-bench.  These long-horizon tasks --- often requiring multi-file patches and hours of professional developer effort --- directly motivate our continuation mechanism.  \citet{prathifkumar2025swebenchverified} raises the important question of benchmark contamination, showing that overlap between SWE-Bench-Verified problems and LLM training data may inflate scores, suggesting that high benchmark numbers may partly reflect memorization rather than genuine problem-solving ability. \citet{horikawa2025agentic} provides an empirical study of how AI coding agents handle refactoring tasks, finding that while agents can plan and execute complex refactorings, they still struggle with cross-file dependency analysis --- a challenge that our sub-agent parallelism partially addresses.

\textbf{Self-improvement and test-time scaling.} \citet{robeyns2025selfimproving} demonstrates a self-improving coding agent that iteratively refines its own scaffolding code, achieving dramatic benchmark improvements through self-generated optimizations. Our self-improvement loop (via \texttt{USER\_PREFS.md}) captures a lighter-weight form of this idea: rather than rewriting its own code, the agent accumulates learned user preferences and project invariants across sessions.  \citet{gao2025trae} introduces the Trae Agent, which applies test-time compute scaling to software engineering, dynamically allocating more inference budget to harder problems.  Our budget-tracking mechanism provides a complementary perspective: rather than scaling compute adaptively, we enforce hard budget ceilings while the Relentless Agent ensures maximum progress within those bounds.

%----------------------------------------------------------------------
\section{Conclusion}
\label{sec:conclusion}
%----------------------------------------------------------------------

We have shown that a simple layered, single-concern architecture can address the practical challenges of deploying LLM agents for real-world software development.  On Terminal Bench~2.0, a benchmark of 89 diverse terminal-based tasks evaluated across 5 trials each, our system achieves a 62.2\% overall pass rate (277/445)---outperforming Claude~Code (58\%) and Cursor Composer~2 (61.7\%) on the same benchmark with the same underlying model (Claude Opus~4.6).  These results are achieved without benchmark-specific optimizations, fine-tuning, or reinforcement learning: the entire agent is roughly 1,900 lines of straightforward Python.

We complement the architecture with a system prompt that encodes engineering discipline directly into the agent's behavior: read before writing, test before fixing, plan before executing, verify before finishing.  Both KISS Sorcar and the KISS Agent Framework are grounded in disciplined engineering practice; these principles are encoded directly into the agent's system prompt, enabling KISS Sorcar to write code that is simple, elegant, maintainable, and bug-free.  These are not novel insights---they are the practices of careful software engineering, translated into instructions that an LLM can follow.  The evaluation suggests that giving a frontier model the time and tools to validate its own output matters more than model-level customization.

In summary, we showed that a simple agent framework, without sophisticated agent technologies such as trajectory compaction and asynchronous multi-agent orchestration, was sufficient to build KISS Sorcar.  By building KISS Sorcar using itself and beating both the Cursor and Claude Code agents, we demonstrated that time-tested software engineering techniques and principles are invaluable for building reliable, practical agent systems.

\section*{An Important Advice} We found that researchers are publishing 
lots of papers on AI, and it is hard to keep track of them or validate their claims. We propose to use the following prompt with KISS Sorcar to identify issues in blogs, papers, and code repositories:

\begin{prompt}
Can you read <<url>>, and thoroughly and precisely check for **wrong assumptions**, **cheating**,
**irreproducibility issues**, **fraud**, **potential for cheating in evaluation**
and **security vulnerabilties**?  
Use the internet search extensively and do not believe what people say--verify it yourself.  
Do not hesitate to download the code and run it to validate the results.
For security vulnerabilities, create a POC and test it.  
Generate an HTML report in PWD/sorcar_reported_frauds/ and open it in the user's default browser.
Thoroughly fact check everything you claim in the report.
\end{prompt}
In the prompt, replace <<url>> with the actual link to a blog, a paper, or a code repository.

\section*{Acknowledgments}

This research is supported in part by gifts from Accenture, Amazon, AMD, Anyscale, Broadcom, Google, IBM, Intel, Intesa Sanpaolo, Lambda, Lightspeed, Mibura, NVIDIA, Samsung SDS, SAP, by the U.S. Department of Energy, Office of Science, Office of Advanced Scientific Computing Research through the X-STACK: Programming Environments for Scientific Computing program (DESC0021982,) and the Defense Advanced Research Projects Agency (DARPA) under Agreement No. HR00112590134.

We would like to thank Marius Momeu for finding critical bugs in the cost computation and in the UI usability, Yogya Mehrotra for testing and fixing bugs in the OpenClaw like third\_party\_agents in KISS Sorcar; Debabrata Dash, Yiwei Hou, Kaiyao Ke, Muxi Lyu, Anoop Mishra, Manish Shetty, Ion Stoica, Shangyin Tan, Hao Wang, and Matei Zaharia for various useful feedback.

Disclaimer: The paper has been edited in part using KISS Sorcar, an AI tool.

%----------------------------------------------------------------------
% References
%----------------------------------------------------------------------


{ \small

\bibliographystyle{plainnat}
\bibliography{kiss_sorcar}
}

\end{document}
